\chapter{Digestion and Circulation}\label{ch:g7:digestion}

\omterm{def:g2:food-energy:food}{Food} must be broken into small molecules that can enter the \omterm{def:g7:digestion:blood}{blood}; \omterm{def:g7:digestion:blood}{blood}
must deliver those molecules (and oxygen) to every \omterm{def:g7:microscope:cell}{cell}. This chapter
links the digestive and \omterm{def:g4:body-systems:circ}{circulatory systems}.

\section{Digestion}

\begin{definition}[Digestion]\label{def:g7:digestion:def}
\emph{Digestion}\index{digestion} is the breakdown of large \omterm{def:g2:food-energy:food}{food} molecules
into smaller soluble molecules that can be absorbed. It includes
mechanical breakdown (chewing, churning) and chemical breakdown (enzymes).
\end{definition}

\begin{definition}[Alimentary canal]\label{def:g7:digestion:canal}
The \emph{alimentary canal}\index{alimentary canal} is the long tube from
\omterm{def:g3:teeth:path}{mouth} to anus: \omterm{def:g3:teeth:path}{mouth}, oesophagus, \omterm{def:g3:teeth:path}{stomach}, small \omterm{def:g3:teeth:path}{intestine}, large
\omterm{def:g3:teeth:path}{intestine}. Associated \omterm{def:g7:tissues:organ}{organs} (liver, pancreas, gall bladder) add juices
but are not part of the tube itself.
\end{definition}

\omimg{digestive-system.jpg}{0.52\linewidth}{Human \omterm{def:g4:body-systems:digestion}{digestive system}:
\omterm{def:g7:tissues:organ}{organs} along the \omterm{def:g7:digestion:canal}{alimentary canal} and accessory glands.}

\begin{omfigure}
\begin{tikzpicture}[font=\footnotesize,
  s/.style={draw, rounded corners, align=center, minimum width=2.0cm,
    minimum height=0.55cm, fill=omDef!12}]
  \node[s] (m) at (0,5.5) {mouth};
  \node[s] (o) at (0,4.4) {oesophagus};
  \node[s, fill=omMeth!15] (st) at (0,3.3) {stomach};
  \node[s, fill=omProp!15] (si) at (0,2.2) {small intestine};
  \node[s, fill=omThm!15] (li) at (0,1.1) {large intestine};
  \node[s] (a) at (0,0) {anus};
  \foreach \a/\b in {m/o,o/st,st/si,si/li,li/a}
    \draw[-Stealth, thick] (\a) -- (\b);
  \node[anchor=west, align=left] at (1.8,3.3)
    {path of food\\(schematic)};
  \node[anchor=west, omExo, align=left] at (1.8,2.2)
    {main absorption\\here};
\end{tikzpicture}

{\small Follow the path: mechanical and chemical \omterm{def:g7:digestion:def}{digestion} prepare
\omterm{def:g4:plants-need:needs}{nutrients} for absorption mainly in the small \omterm{def:g3:teeth:path}{intestine}.}
\end{omfigure}

\begin{proposition}[Key jobs along the tube]\label{prop:g7:digestion:jobs}
\begin{itemize}
  \item \omterm{def:g3:teeth:path}{Mouth}: chew; saliva starts starch \omterm{def:g7:digestion:def}{digestion}.
  \item \omterm{def:g3:teeth:path}{Stomach}: acid and enzymes begin \omterm{def:g6:nutrition:groups}{protein} \omterm{def:g7:digestion:def}{digestion}; churning.
  \item Small \omterm{def:g3:teeth:path}{intestine}: most chemical \omterm{def:g7:digestion:def}{digestion} finished; products
        absorbed into \omterm{def:g7:digestion:blood}{blood} (and \omterm{def:g6:nutrition:groups}{fats} into lymph).
  \item Large \omterm{def:g3:teeth:path}{intestine}: \omterm{def:g4:plants-need:needs}{water} reabsorbed; faeces formed.
\end{itemize}
\end{proposition}
\admitted

\begin{definition}[Absorption and villi]\label{def:g7:digestion:villi}
\emph{Absorption}\index{absorption (digestion)} is the movement of
digested \omterm{def:g4:plants-need:needs}{nutrients} into the \omterm{def:g7:digestion:blood}{blood} or lymph. The small \omterm{def:g3:teeth:path}{intestine} lining has
finger-like \emph{villi}\index{villus} that increase surface area.
\end{definition}

\section{Circulation}

\begin{definition}[Blood and vessels]\label{def:g7:digestion:blood}
\emph{Blood}\index{blood} carries oxygen, \omterm{def:g4:plants-need:needs}{nutrients}, hormones, wastes and
defences. \emph{Arteries}\index{artery} carry blood away from the \omterm{def:g7:digestion:heart}{heart};
\emph{veins}\index{vein} return blood to the \omterm{def:g7:digestion:heart}{heart}; \emph{capillaries}
\index{capillary} are tiny exchange vessels in \omterm{def:g7:tissues:tissue}{tissues}.
\end{definition}

\begin{definition}[Heart]\label{def:g7:digestion:heart}
The \emph{heart}\index{heart} is a muscular pump with four chambers in
\omterm{def:g2:animal-groups:five}{mammals}: two atria and two ventricles. It drives a double circuit: one to
the \omterm{def:g5:circulation:breathing}{lungs} (pick up oxygen) and one to the \omterm{def:g1:my-body:body}{body}.
\end{definition}

\begin{omfigure}
\begin{tikzpicture}[font=\footnotesize, scale=0.95]
  \draw[thick, omExo, fill=omExo!10] (0,0) circle (1.1);
  \node at (0,0) {heart};
  \draw[thick, omThm, fill=omThm!12] (4.5,1.6) circle (0.9);
  \node at (4.5,1.6) {lungs};
  \draw[thick, omDef, fill=omDef!12] (4.5,-1.6) circle (0.9);
  \node at (4.5,-1.6) {body};
  \draw[-Stealth, very thick, omThm] (0.9,0.5) -- (3.6,1.4);
  \draw[-Stealth, very thick, omProp] (3.6,1.8) -- (0.9,0.7);
  \draw[-Stealth, very thick, omDef] (0.9,-0.5) -- (3.6,-1.4);
  \draw[-Stealth, very thick, omMeth] (3.6,-1.8) -- (0.9,-0.7);
  \node[font=\scriptsize] at (2.2,1.5) {to lungs};
  \node[font=\scriptsize] at (2.2,0.9) {from lungs};
  \node[font=\scriptsize] at (2.2,-0.9) {to body};
  \node[font=\scriptsize] at (2.2,-1.5) {from body};
\end{tikzpicture}

{\small Double circulation idea: \omterm{def:g7:digestion:heart}{heart} $\leftrightarrow$ \omterm{def:g5:circulation:breathing}{lungs} and \omterm{def:g7:digestion:heart}{heart}
$\leftrightarrow$ \omterm{def:g1:my-body:body}{body}.}
\end{omfigure}

\begin{proposition}[Link to digestion]\label{prop:g7:digestion:link}
Absorbed glucose and amino acids enter \omterm{def:g7:digestion:blood}{blood} at the small \omterm{def:g3:teeth:path}{intestine} and
travel via the liver pathway into general circulation. \omterm{def:g7:microscope:cell}{Cells} use them for
\omterm{def:g6:nutrition:def}{energy} and building; the \omterm{def:g7:digestion:heart}{heart} keeps the delivery system moving.
\end{proposition}
\admitted

\begin{example}[Pulse]\label{ex:g7:digestion:pulse}
A pulse is the pressure wave in arteries from heartbeats. Resting pulse
and exercise pulse differ because \omterm{def:g3:bones:muscle}{muscles} need more oxygen and glucose
when active.
\end{example}

\section{Exercises}

\begin{exercise}[$\star$]\label{exo:g7:digestion:1}
Define \omterm{def:g7:digestion:def}{digestion}. Why is it needed before absorption?
\end{exercise}

\begin{exercise}[$\star$]\label{exo:g7:digestion:2}
List the order of \omterm{def:g7:tissues:organ}{organs} \omterm{def:g2:food-energy:food}{food} passes through from \omterm{def:g3:teeth:path}{mouth} to anus.
\end{exercise}

\begin{exercise}[$\star$]\label{exo:g7:digestion:3}
Where does most absorption of digested \omterm{def:g2:food-energy:food}{food} occur? Name the structures
that increase surface area there.
\end{exercise}

\begin{exercise}[$\star$]\label{exo:g7:digestion:4}
Distinguish \omterm{def:g7:digestion:blood}{artery}, \omterm{def:g7:digestion:blood}{vein} and \omterm{def:g7:digestion:blood}{capillary}.
\end{exercise}

\begin{exercise}[$\star$]\label{exo:g7:digestion:5}
Why is the mammalian circulation called a double circulation?
\end{exercise}

